\paragraph{Tres niveles de evaluación.} Se distinguen: (i) los \emph{criterios de aceptación originales}, fijados antes de simular con un umbral para cada uno (sección~\ref{sec:calib}); con el terreno anterior ninguna simulación los cumplía todos; la candidata a la prueba reservada (z02) cumple tres de cuatro (no cumple el de perímetro), mientras que los ajustes posteriores de casas porosas (z26) y de la hipótesis integrada con lluvia supuesta (z32) cumplen los cuatro; (ii) el \emph{índice exploratorio de compatibilidad} (IEC), definido después, que resume los mismos cuatro criterios en una escala continua para comparar hipótesis; y (iii) la \emph{validación independiente} de los ajustes posteriores, no probada al cierre, que exige datos no usados para formular las hipótesis, como imágenes satelitales posteriores al evento, levantamientos oficiales de depósitos o marcas de nivel.

\paragraph{Índice exploratorio de compatibilidad.} Para comparar hipótesis se resume cada simulación en un \emph{índice exploratorio de compatibilidad} (IEC) con lo observado, en una escala de 0 a 100: el promedio de cuatro componentes, cada uno entre 0 y 1, derivados de los criterios originales fijados antes de simular (sección~\ref{sec:calib}); el índice se definió después y la referencia de perímetro se actualizó con el polígono oficial. Son la fracción de puntos de las calles con barro informado por las fuentes que el modelo moja ($R$); la fracción de los cuatro recintos con afectación informada alcanzados ($R_\mathrm{sit}$); el alcance, igual a 1 si el barro llega a menos de \SI{200}{m} de Padre Hurtado y decreciente en forma lineal hasta 0 a \SI{1200}{m}; y $F_\mathrm{per}$. El IEC no es una probabilidad ni una medida de precisión predictiva: resume qué tan bien una simulación reproduce un conjunto de observaciones cualitativas, y diferencias de dos o tres puntos no deben interpretarse como evidencia concluyente de superioridad entre hipótesis. Las profundidades, velocidades y presiones por calle y edificio se extraen como el máximo en un radio de \SI{40}{m} de cada punto de interés.
